\chapter*{Nouveautés}
\addstarredchapter{\textbf{Nouveautés}}

Voici la liste des derniers correctifs et des dernières implémentations de fonctionnalité dans l'application \cwv.

\begin{table}[H]
    \centering
    \begin{tabular}{>{\raggedright}p{3.25cm}>{\raggedright}p{12cm}>{\raggedright}p{2cm}}
    \hline
    \textbf{Type/Tag} & \textbf{Description} & \textbf{Section}\tabularnewline
    \hline
    \newfeature & Séparation du backend et du frontend dans deux dépôts distincts reliés par des API. Le backend constitue la première version du Jumeau Hydrologique : Hydrological Twin Alpha Series (HTAS) & Sec. \ref{sec-generalites} \tabularnewline
    \arrayrulecolor{lightgray} \hline
    \newfeature & Automatisation de la configuration : seul le répertoire de sorties CaWaQS est nécessaire & Sec. \ref{sec_config_manuelle} \tabularnewline
    \arrayrulecolor{lightgray} \hline
    \newfeature & Extraction par masque. Découpage des résultats sur une sous-zone à partir d'un shapefile & Sec. \ref{sec-extract-area} \tabularnewline
    \arrayrulecolor{lightgray} \hline
    \end{tabular}
    \caption{Dernières implémentations dans l'application}
\end{table}




\textbf{Derniers tags stables :} \\

Tag CWV\_alpha\_3.72 : \href{https://gitlab.com/cawaqs/gviz/cawaqsviz/-/tree/CWV_alpha_3.72?ref_type=tags}{Lien} ~ (05 juin 2026) : tag stable de \cwv{} couplé à HTAS.
